\chapter{Introducción}
\label{chap:intro}

\ph{Este capítulo presenta el trabajo. Da al lector el contexto mínimo para entender
qué se aborda, por qué es relevante, qué se pretende conseguir y qué queda fuera.
No se entra todavía en detalles técnicos: eso pertenece al estado del arte y al diseño.
Fija aquí el registro que mantendrás en toda la memoria: redacción impersonal
(``se ha desarrollado'') o primera persona del plural (``hemos desarrollado''),
pero consistente de principio a fin, y prosa académica sin coloquialismos
(consulta la Guía de estilo incluida al inicio del documento en modo borrador).}

\section{Contexto}
\label{sec:contexto}

\ph{Describe el dominio en el que se enmarca el TFG y por qué es relevante.
Aporta el trasfondo necesario para que un lector ajeno al tema entienda
el resto del documento. Un lector tras leer esta sección debería poder responder a:
``¿de qué va este trabajo y en qué entorno se sitúa?''.}

\ph{Sugerencia de contenido:}
\begin{itemize}
  \item \ph{Sector o disciplina (industria, investigación, educación\dots).}
  \item \ph{Tendencias o cambios recientes que motivan el problema.}
  \item \ph{Conceptos básicos imprescindibles para entender el resto.}
\end{itemize}

\section{Problema}
\label{sec:problema}

\ph{Plantea el problema concreto que el TFG resuelve. Conviene formularlo en términos
de limitaciones reales del estado actual: qué no se puede hacer hoy o qué se hace mal,
y qué consecuencias tiene.}

\ph{Estructura recomendada (1--3 problemas claramente delimitados):}

\begin{enumerate}
  \item \textbf{\ph{Limitación 1:}} \ph{descripción del problema y de su impacto.}
  \item \textbf{\ph{Limitación 2:}} \ph{descripción del problema y de su impacto.}
  \item \textbf{\ph{Limitación 3:}} \ph{descripción del problema y de su impacto.}
\end{enumerate}

\section{Motivación}
\label{sec:motivacion}

\ph{Justifica por qué merece la pena resolver el problema y por qué con este enfoque.
Conecta el problema (sección anterior) con la solución que se va a proponer
(sin entrar todavía en cómo se implementa).}

\ph{Puntos típicos:}
\begin{itemize}
  \item \ph{Por qué el enfoque elegido es adecuado.}
  \item \ph{Qué aporta respecto a soluciones existentes.}
  \item \ph{Encaje del trabajo en un contexto académico, profesional o de investigación
        (grupo de investigación, asignatura, empresa, publicación\dots).}
\end{itemize}

\section{Objetivos}
\label{sec:objetivos}

\ph{Define el objetivo general del trabajo en una o dos frases, y a continuación
los objetivos específicos en forma de lista. Formula cada objetivo específico
siguiendo los criterios \textbf{SMART}: \textbf{S}pecific (específico: una sola
cosa, sin ambigüedad), \textbf{M}easurable (medible: con una métrica o criterio
de éxito verificable), \textbf{A}chievable (alcanzable: realista con los recursos
y el plazo de un TFG), \textbf{R}elevant (relevante: conectado con una de las
limitaciones de la Sección~\ref{sec:problema}) y \textbf{T}ime-bound (acotado en
el tiempo: realizable dentro del calendario del Capítulo~\ref{chap:planificacion}).}

\ph{Usa verbos de resultado, no de actividad: evita ``estudiar'', ``aprender'' o
``profundizar en'' (no se puede verificar cuándo terminan) y prefiere
``implementar X que soporte Y'', ``alcanzar una cobertura del X\,\%'', ``reducir
el tiempo de Z por debajo de N segundos''. Ejemplo de objetivo bien formulado:
``Implementar un módulo de búsqueda que resuelva consultas sobre 10\,000
registros en menos de 200~ms, dando respuesta a la Limitación~1''.}

\ph{Objetivo general: ``Desarrollar/implementar/diseñar \dots con el fin de \dots''.}

\begin{enumerate}[label=\textbf{O\arabic*.}]
  \item \textbf{\ph{Objetivo específico 1:}} \ph{descripción concreta y medible.}
  \item \textbf{\ph{Objetivo específico 2:}} \ph{descripción concreta y medible.}
  \item \textbf{\ph{Objetivo específico 3:}} \ph{descripción concreta y medible.}
  \item \textbf{\ph{Objetivo específico 4:}} \ph{descripción concreta y medible.}
  \item \textbf{\ph{Objetivo específico 5 (opcional):}} \ph{descripción concreta y medible.}
\end{enumerate}

\ph{Resume en la tabla siguiente el criterio de éxito de cada objetivo, cómo se
verificará y en qué capítulo. Esta tabla es la columna vertebral de la memoria:
el Capítulo~\ref{chap:resultados} dará el veredicto sobre \textbf{todos} los
criterios, recapitulando con \textbackslash ref la evidencia de aquellos que se
verifican en capítulos anteriores, y el Capítulo~\ref{chap:conclusiones}
sintetizará el grado de cumplimiento final. Si el criterio es cuantitativo,
fíjalo también como RNF en el Capítulo~\ref{chap:requisitos} con el
\textbf{mismo} umbral: un único umbral, dos vistas. Conviene rellenar la
columna «Cap.» con \textbackslash ref (p.\,ej.\ Capítulo~\ref{chap:pruebas})
para que la numeración se mantenga sola.}

\begin{longtable}{@{}P{1cm} P{6.2cm} P{4.6cm} P{1.7cm}@{}}
  \toprule
  \textbf{Obj.} & \textbf{Criterio de éxito (medible)} & \textbf{Cómo se verifica} & \textbf{Cap.} \\
  \midrule
  \endhead
  O1 & \ph{p.\,ej.\ el sistema cubre todos los RF de prioridad alta} & \ph{matriz de trazabilidad y suite de pruebas} & \ph{4, 7} \\[0.3em]
  O2 & \ph{p.\,ej.\ tiempo de respuesta inferior a 200~ms en consultas típicas} & \ph{medición sobre el banco de pruebas} & \ph{8} \\[0.3em]
  O3 & \ph{criterio medible con métrica y umbral.} & \ph{evidencia que lo demuestra.} & \ph{X} \\[0.3em]
  O4 & \ph{criterio medible con métrica y umbral.} & \ph{evidencia que lo demuestra.} & \ph{X} \\[0.3em]
  O5 & \ph{criterio medible con métrica y umbral.} & \ph{evidencia que lo demuestra.} & \ph{X} \\
  \bottomrule
  \caption{Criterios de éxito de los objetivos y capítulos donde se verifican.}
  \label{tab:criterios-exito}
\end{longtable}

\section{Alcance}
\label{sec:alcance}

\ph{Delimita lo que el trabajo \textbf{sí} cubre y, sobre todo, lo que \textbf{no} cubre.
Es la sección que protege al autor de expectativas desmedidas durante la defensa.}

\ph{Una posible estructura:}

\ph{Lo que sí se aborda:}
\begin{itemize}
  \item \ph{Funcionalidad / dominio / tipo de entradas que el sistema acepta.}
  \item \ph{Restricciones cuantitativas (rangos, tamaños, plataformas\dots).}
\end{itemize}

\ph{Quedan explícitamente fuera del alcance:}
\begin{itemize}
  \item \ph{Funcionalidad excluida 1 (con una breve justificación).}
  \item \ph{Funcionalidad excluida 2 (con una breve justificación).}
  \item \ph{Funcionalidad excluida 3 (con una breve justificación).}
\end{itemize}

\section{Estructura de la memoria}
\label{sec:estructura_memoria}

\ph{Cierra la introducción con un mapa de lectura: una o dos frases por capítulo.
No describas solo el contenido (``se presentan los requisitos''), sino el papel
de cada capítulo en el argumento de la memoria (``se establece qué debe hacer el
sistema, base contractual del diseño'').}

El resto de la memoria se organiza como sigue:

\begin{itemize}
  \item El \textbf{Capítulo~\ref{chap:estado_arte}} \ph{revisa los fundamentos
        teóricos y los trabajos relacionados, identificando el hueco que este
        trabajo cubre.}
  \item El \textbf{Capítulo~\ref{chap:planificacion}} \ph{describe la metodología,
        las fases, los riesgos y el presupuesto con los que se ha gestionado el
        proyecto.}
  \item El \textbf{Capítulo~\ref{chap:requisitos}} \ph{especifica qué debe hacer
        el sistema: requisitos, casos de uso y su trazabilidad, base contractual
        del resto del trabajo.}
  \item El \textbf{Capítulo~\ref{chap:diseno}} \ph{presenta cómo se estructura la
        solución: arquitectura, modelos y decisiones de diseño que responden a
        los requisitos.}
  \item El \textbf{Capítulo~\ref{chap:implementacion}} \ph{detalla cómo se
        materializa el diseño: tecnologías elegidas y módulos construidos.}
  \item El \textbf{Capítulo~\ref{chap:pruebas}} \ph{muestra cómo se ha verificado
        que el sistema cumple los requisitos y con qué estrategia de pruebas.}
  \item El \textbf{Capítulo~\ref{chap:resultados}} \ph{contrasta los resultados
        obtenidos con los criterios de éxito de la Tabla~\ref{tab:criterios-exito}
        y los discute frente a los trabajos relacionados.}
  \item El \textbf{Capítulo~\ref{chap:conclusiones}} \ph{da el veredicto sobre el
        cumplimiento de los objetivos, reconoce las limitaciones y propone líneas
        de trabajo futuro.}
\end{itemize}

% --- Transición al siguiente capítulo (último párrafo del roadmap) ---
\ph{Remata el roadmap con un párrafo puente de 2--4 líneas que lance al lector al
siguiente capítulo. Por ejemplo: «Definidos el problema, los objetivos y el
alcance, el Capítulo~\ref{chap:estado_arte} revisa los fundamentos y las
soluciones existentes sobre los que se construye la propuesta».}
